\documentclass[12pt,a4paper]{article}
\usepackage[T1]{fontenc}
\usepackage{amsmath,amssymb,mathrsfs,tikz,times,pifont}
\usepackage{enumitem}
\usepackage{multicol}
\usepackage{lmodern}
\usetikzlibrary{trees}
\newcommand\circitem[1]{%
\tikz[baseline=(char.base)]{
\node[circle,draw=gray, fill=red!55,
minimum size=1.2em,inner sep=0] (char) {#1};}}
\newcommand\boxitem[1]{%
\tikz[baseline=(char.base)]{
\node[fill=cyan,
minimum size=1.2em,inner sep=0] (char) {#1};}}
\setlist[enumerate,1]{label=\protect\circitem{\arabic*}}
\setlist[enumerate,2]{label=\protect\boxitem{\alph*}}
%%%::::::by chnini ameur :::::::%%%
\everymath{\displaystyle}
\usepackage[left=1cm,right=1cm,top=1cm,bottom=1.7cm]{geometry}
\usepackage[colorlinks=true, linkcolor=blue, urlcolor=blue, citecolor=blue]{hyperref}
\usepackage{array,multirow}
\usepackage[most]{tcolorbox}
\usepackage{varwidth}
\usepackage{float} %pour utiliser l'option [H] qui force l'image à apparaître exactement à l'endroit où elle est placée dans le code.
\tcbuselibrary{skins,hooks}
\usetikzlibrary{patterns}
%%%::::::by chnini ameur :::::::%%%
\newtcolorbox{exa}[2][]{enhanced,breakable,before skip=2mm,after skip=5mm,
colback=yellow!20!white,colframe=black!20!blue,boxrule=0.5mm,
attach boxed title to top left ={xshift=0.6cm,yshift*=1mm-\tcboxedtitleheight},
fonttitle=\bfseries,
title={#2},#1,
% varwidth boxed title*=-3cm,
boxed title style={frame code={
\path[fill=tcbcolback!30!black]
([yshift=-1mm,xshift=-1mm]frame.north west)
arc[start angle=0,end angle=180,radius=1mm]
([yshift=-1mm,xshift=1mm]frame.north east)
arc[start angle=180,end angle=0,radius=1mm];
\path[left color=tcbcolback!60!black,right color = tcbcolback!60!black,
middle color = tcbcolback!80!black]
([xshift=-2mm]frame.north west) -- ([xshift=2mm]frame.north east)
[rounded corners=1mm]-- ([xshift=1mm,yshift=-1mm]frame.north east)
-- (frame.south east) -- (frame.south west)
-- ([xshift=-1mm,yshift=-1mm]frame.north west)
[sharp corners]-- cycle;
},interior engine=empty,
},interior style={top color=yellow!5}}
%%%%%%%%%%%%%%%%%%%%%%%
\usepackage{fancyhdr}
\usepackage{eso-pic}         % Pour ajouter des éléments en arrière-plan
% Commande pour ajouter du texte en arrière-plan
\usepackage{tkz-tab}
\usepackage{tikz}
\usepackage[siunitx, european, RPvoltages]{circuitikz}
\AddToShipoutPicture{
    \AtTextCenter{%
        \makebox[0pt]{\rotatebox{80}{\textcolor[gray]{0.7}{\fontsize{5cm}{5cm}\selectfont PGB}}}
    }
}
\usepackage{lastpage}
\fancyhf{}
\pagestyle{fancy}
\renewcommand{\footrulewidth}{1pt}
\renewcommand{\headrulewidth}{0pt}
\renewcommand{\footruleskip}{10pt}
\fancyfoot[R]{
\color{blue}\ding{45}\ \textbf{2026}
}
\fancyfoot[L]{
\color{blue}\ding{45}\ \textbf{Prof:M. BA}
}
\cfoot{\bf
\thepage /
\pageref{LastPage}}
% Création du compteur pour les exercices
\newcounter{exercice}
\renewcommand{\theexercice}{\arabic{exercice}}  % Définit l'affichage du compteur en chiffres arabes

% Définir la commande \exo
\newcommand{\exo}{\refstepcounter{exercice}\textbf{Exercice \theexercice} }
\begin{document}
\renewcommand{\arraystretch}{1.5}
\renewcommand{\arrayrulewidth}{1.2pt}
\begin{tikzpicture}[overlay,remember picture]
    \node[draw=blue,line width=1.2pt,fill=purple,text=blue,inner sep=3mm,rounded corners,pattern=dots]at ([yshift=-2.5cm]current page.north) {\begingroup\setlength{\fboxsep}{0pt}\colorbox{white}{\begin{tabular}{|*1{>{\centering \arraybackslash}p{0.28\textwidth}} |*2{>{\centering \arraybackslash}p{0.2\textwidth}|} *1{>{\centering \arraybackslash}p{0.19\textwidth}|} }
                \hline
                \multicolumn{3}{|c|}{$\diamond$$\diamond$$\diamond$\ \textbf{Lycée de Dindéfélo}\ $\diamond$$\diamond$$\diamond$ } & \textbf{A.S. : 2026/2027}                                              \\ \hline
                \textbf{Matière: Mathématiques}                                                                                    & \textbf{Niveau : 3}\textbf{M1} & \textbf{Date: 18/10/2026} & \textbf{} \\ \hline
                \multicolumn{4}{|c|}{\parbox[c]{10cm}{\begin{center}
                                                                  \textbf{{\Large\sffamily Racine Carré}}
                                                              \end{center}}}                                                                                                        \\ \hline
            \end{tabular}}\endgroup};
\end{tikzpicture}
\vspace{3cm}
\begin{multicols}{2}
\setlength{\columnseprule}{0.1mm} % La largeur de la ligne verticale entre les colonnes
\fbox{\textbf{\exo}}
Écrire sous la forme $a\sqrt{b}$ où $a$ et $b$ sont des nombres entiers, $b$ étant le plus petit possible :
\begin{enumerate}
        \item $A=\sqrt{50}$
        \item $B=\sqrt{72}$
        \item $C=\sqrt{50} + \sqrt{72}$
        \item $D=2\sqrt{3} + \sqrt{75} - 6\sqrt{27}$
        \item $E=2\sqrt{3} \times \sqrt{6}$
        \item $F=\sqrt{8} \times \sqrt{50} \times \sqrt{18}$
\end{enumerate}

\fbox{\textbf{\exo}}
Écrire les nombres suivants sous la forme $p + q\sqrt{7}$ où $p$ et $q$ sont des entiers relatifs :
\begin{enumerate}
\item $A=\sqrt{49} + \sqrt{28} + \sqrt{63}$
\item $B=(2\sqrt{7} + 1)^2 - (\sqrt{3} - 1)(\sqrt{3} + 1)$
\item $C=6\sqrt{28} + 10\sqrt{7} - 8\sqrt{63}$
\end{enumerate}

\fbox{\textbf{\exo}}
On considère les nombres $D$ et $E$ suivants : $D = (2\sqrt{3} + 1) \times (2\sqrt{3} - 1)$ et\\ $E = 8\sqrt{5} - \sqrt{20} - 2\sqrt{45}$.

En indiquant le détail des calculs, écrire $D$ et $E$ sous la forme de nombres entiers.

\fbox{\textbf{\exo}}
\begin{enumerate}
\item Écrire $\sqrt{5} \times \sqrt{125}$ sous la forme d'un nombre entier.
\item Écrire $(\sqrt{5} \times \sqrt{125}) \times 2$ sous la forme $a\sqrt{5}$ où $a$ est un entier.
\end{enumerate}

\fbox{\textbf{\exo}}

\begin{enumerate}
\item Écrire $A$, $B$ et $C$ sous la forme $a\sqrt{3}$ où $a$ est un entier.
\begin{enumerate}
\item $A= \sqrt{12} + 5\sqrt{75} - 2\sqrt{27}$
\item $B= (5 + \sqrt{3})^2 - (2\sqrt{7})^2$
\item $C= \sqrt{12} - \sqrt{75} - 2\sqrt{27}$
\end{enumerate}
\item On donne : $E = 3\sqrt{5} - 2\sqrt{11}$ et $F = 3\sqrt{5} + 2\sqrt{11}$. Écrire et calculer le produit des nombres $E$ et $F$.
\end{enumerate}

\fbox{\textbf{\exo}} % Combiné d'après les captures
On pose $A = \sqrt{48} + \sqrt{20}$ et $B = \sqrt{108} - \sqrt{45}$.
\begin{enumerate}
\item Montrer que :
\begin{enumerate}
\item $A$ s'écrit sous la forme $a\sqrt{3} + b\sqrt{5}$
\item $B$ s'écrit sous la forme $c\sqrt{3} + d\sqrt{5}$ où $a, b, c, d$ sont des entiers relatifs.
\end{enumerate}
\item Montrer que le produit $AB$ est un nombre entier.
\end{enumerate}

\fbox{\textbf{\exo}}

On pose : $A = \sqrt{27} + 1$ ; $B = 2\sqrt{3} - 5$. 

Écrire sous la forme $a\sqrt{3} + b$, où $a$ et $b$ sont deux entiers relatifs, les nombres suivants : A - B ; $A^2$ et $B^2$.

\fbox{\textbf{\exo}}

On donne les nombres $D = 5 - 3\sqrt{2}$ et 

$E = 4 + 5\sqrt{2}$.

Calculer D - E ; D * E. 

Donner les résultats sous la forme $a + b\sqrt{2}$ où $a$ et $b$ sont des nombres entiers relatifs.

\fbox{\textbf{\exo}}

\begin{enumerate}
\item Soit $C = \sqrt{500} + 3\sqrt{5} - 3\sqrt{45}$

Écrire $C$ sous la forme $a\sqrt{b}$ où $a$ et $b$ sont des entiers, $b$ étant le plus petit possible.
\item Soit $D = (5 - 2\sqrt{6})(5 + 2\sqrt{6})$. 

Exprimer $D$ sous forme d'un nombre entier.
\end{enumerate}

\fbox{\textbf{\exo}}

On pose $B = (5 + \sqrt{3})(5 - \sqrt{3}) - 8\sqrt{5}(\sqrt{5} - 1)$. 

Écrire $B$ sous forme $a + b\sqrt{5}$ (avec $a$ et $b$ étant des nombres relatifs).

\fbox{\textbf{\exo}}

On pose : $a = \sqrt{3}(1 + \sqrt{6})$ et $b = 3 - \sqrt{6}$

\begin{enumerate}
\item Calculer $a^2$ ; $b^2$ et $a^2 + b^2$
\item Montrer que $a^2 + b^2$ est un nombre entier.
\item Si $a$ et $b$ sont les longueurs des côtés de l'angle droit dans un triangle rectangle, quelle est la longueur de l'hypoténuse ?
\end{enumerate}

\fbox{\textbf{\exo}}

Simplifier les expressions suivantes :
\begin{enumerate}
\item $A=\dfrac{\sqrt{8} + 2\sqrt{32}}{5\sqrt{2} - \sqrt{50}} $
\item $B=\sqrt{2 - \sqrt{3}} \times \sqrt{2 + \sqrt{3}} $
\item $C=\frac{\sqrt{2} + 1}{\sqrt{2} - 1} + \frac{2\sqrt{2} + 1}{\sqrt{2} + 1}$
\end{enumerate}
 
\fbox{\textbf{\exo}}

On donne : $a = \frac{2 - \sqrt{3}}{5 + \sqrt{3}}$

$b = 3\sqrt{18} + \sqrt{128} - \sqrt{338} \qquad c = \sqrt{2} - 3$

\begin{enumerate}
\item Rendre rationnel le dénominateur de $a$ puis simplifier $b$.
\item Calculer $c^2$. 

En déduire que $p = \frac{6 - \sqrt{8}}{3\sqrt{5} - 6\sqrt{2}}$ est un rationnel que l'on déterminera.
\end{enumerate}

\fbox{\textbf{\exo}}Écrire le plus simplement possible :
\begin{enumerate}
\item $A = \frac{\sqrt{3 - \sqrt{3}} \times \sqrt{3 + \sqrt{3}}}{\sqrt{6}}$
\item $B = \frac{-3(\sqrt{5} - \sqrt{2})}{\sqrt{45} - \sqrt{18}}$
\item $C = \frac{(\sqrt{3} - \sqrt{27} + \sqrt{12})}{\sqrt{54}}$
\item $D = \frac{2 - \sqrt{5}}{2 + \sqrt{5}} - \frac{2 + \sqrt{5}}{2 - \sqrt{5}}$
\item $E = 2 \sqrt{37 - \sqrt{\frac{21}{25} + \frac{1}{25} \times \sqrt{6 + \sqrt{103 - 2\sqrt{\frac{9}{4}}}}}}$
\end{enumerate}

\fbox{\textbf{\exo}}
Rendre rationnel le dénominateur de chacun des nombres suivants :
\begin{multicols}{2}
\setlength{\columnseprule}{0.1mm} % La largeur de la ligne verticale entre les colonnes
\begin{enumerate}
\item $A = \frac{1}{\sqrt{3}}$
\item $B = \frac{\sqrt{2}}{\sqrt{7}}$
\item $C = \frac{1}{\sqrt{11}} - \frac{5}{\sqrt{8}}$
\item $D = \frac{2}{\sqrt{2} + 5}$
\item $E = \frac{1}{\sqrt{3} - \sqrt{5}}$
\item $F = \frac{\sqrt{7} + 1}{3 - \sqrt{2}}$
\item $G = \frac{\sqrt{5}}{\sqrt{3} + 2}$
\item $H = \frac{1}{\sqrt{5} + \sqrt{6}}$
\item $I = \frac{1 + \sqrt{2}}{5 - 3\sqrt{2}}$
\end{enumerate}
\end{multicols}

\fbox{\textbf{\exo}}
\begin{enumerate}
\item Compare les nombres réels suivants 

$\sqrt{5} \text{ et } 2,3\sqrt{2} \qquad\qquad -\sqrt{7} \text{ et } 2\sqrt{7} $ 

$ 5 - \sqrt{2} \text{ et } 3 - \sqrt{2} \qquad\qquad 9 \text{ et } 4\sqrt{5}$

$3\sqrt{2} \text{ et } 2\sqrt{3} \qquad\qquad 2\sqrt{7} \text{ et } -7\sqrt{2}$ 

$\sqrt{-1+\sqrt{2}} \text{ et } \sqrt{1+\sqrt{2}}$
\item Écrire les nombres ci-dessous sans le symbole de la valeur absolue.

$
|9 - 4\sqrt{5}| \qquad\qquad |3\sqrt{2} - 2\sqrt{3}|$ 

$|-2\sqrt{7} - 7\sqrt{2}| \qquad\qquad \left|\sqrt{-1+\sqrt{2}} - \sqrt{1+\sqrt{2}}\right|
$
\end{enumerate}

\fbox{\textbf{\exo}}

On donne : $A = \sqrt{5} + 3$ et $B = \sqrt{5} - 3$

\begin{enumerate}
\item Calculer : $A^2$ ; $B^2$ et $A \times B$ et $A/B$
\item Simplifier $c = \frac{A}{B} + \frac{B}{A}$
\end{enumerate}

\fbox{\textbf{\exo}}
Soient les réels suivants : $a = \frac{\sqrt{7} - 3}{\sqrt{2}}$ et $b = \frac{\sqrt{7} + 3}{-\sqrt{2}}$
\begin{enumerate}
\item Calculer le produit $a \times b$. Que peut-on en déduire pour les réels $a$ et $b$ ?
\item Calculer et comparer les réels $a^2$ et $\frac{a}{b}$
\item Peut-on prévoir ce résultat ?
\end{enumerate}

\fbox{\textbf{\exo}}
On donne : $A = \frac{\sqrt{5} + 1}{2}$ et $2,236 < \sqrt{5} < 2,237$
\begin{enumerate}
\item Écrire l'inverse de $A$ en rendant rationnel le dénominateur.
\item Comparer $A$ et $A^2 - 1$
\item En déduire que $A^2 = A + 1$
\item Donner un encadrement de l'inverse de $A$ et un encadrement du carré de $A$ par deux décimaux consécutifs d'ordre 2.
\end{enumerate}
\end{multicols}

\end{document}